\section{Conclusiones}

El presente trabajo contribuyó al equipo, principalmente, a mejorar el entendimiento en el desarrollo de videojuegos
de tipo MMO, a la vez que permitió profundizar en conocimientos de desarrollo de software, enseñó el verdadero valor
de la gestión y planificación en proyectos de software donde interactúan varios componentes, y la importancia de la 
comunicación entre los miembros del equipo.
\vspace{0.25cm}

En 1986, Fred Brooks en su \textit{paper} \textit{No silver bullet} hablaba, por primera vez, de la complejidad inherente 
al desarrollo de software. Este proyecto fue una prueba más de esa complejidad de la que hablaba Brooks. El equipo tuvo que lidiar con 
diferentes módulos y tecnologías logrando que todos ellos convivan de forma cohesiva. Abordar esa complejidad
fue el resultado de un gran sacrificio y dedicación de un
equipo humano, que a pesar de las vicisitudes y contratiempos, llegó a un resultado final satisfactorio, y logró cumplir con 
los objetivos planteados al inicio del proyecto.
\vspace{0.25cm}

El resultado de este proyecto fue \textbf{Feed The Realm}, una plataforma \textit{MMO} donde los jugadores pueden explorar
mundos creados por la comunidad y donde cualquier usuario puede diseñar, publicar y monetizar sus propias experiencias mediante
la venta de cosméticos adquiribles con gemas. La flexibilidad de las herramientas desarrolladas permite crear mundos de distinta
escala y complejidad, fomentando la diversidad de contenido y la creatividad de los creadores.

Asimismo, el \textit{feedback} obtenido durante las etapas de prueba permitió mejorar distintos aspectos de la experiencia de 
juego y validar decisiones de diseño. 
En paralelo, se desarrolló una infraestructura escalable con despliegue y descubrimiento automático de servidores,
un conjunto de servicios centrales para la gestión de usuarios, mundos, \textit{assets} y pagos, y herramientas de CI/CD y
monitoreo que facilitan la operación y el mantenimiento de la plataforma.

